%%% CSE 4316 Project Charter - WORKED EXAMPLE (Computer Vision / ML)
%%% Structure follows ISO/IEC/IEEE 16326:2019.
\documentclass[11pt,letterpaper]{article}
\usepackage[margin=1.0in]{geometry}
\usepackage[T1]{fontenc}\usepackage[bitstream-charter]{mathdesign}\usepackage[latin1]{inputenc}
\usepackage{amsmath}\usepackage{xcolor}\usepackage{cite}\usepackage{hyphenat}
\usepackage{graphicx}\usepackage{float}\usepackage{subfigure}\usepackage{sectsty}
\usepackage[compact]{titlesec}\usepackage[tablegrid]{vhistory}\usepackage{longtable}
\allsectionsfont{\color{accentcolor}\scshape\selectfont}
\definecolor{accentcolor}{rgb}{0.0,0.0,0.5}
\newcommand{\teamname}{Team Perceptron}
\newcommand{\productname}{DigitDesk Handwritten Form Digitization Service}
\newcommand{\coursename}{CSE 4316: Senior Design I}
\newcommand{\semester}{Fall 2026}
\newcommand{\docname}{Project Charter}
\newcommand{\department}{Department of Computer Science \& Engineering}
\newcommand{\university}{The University of Texas at Arlington}
\newcommand{\authors}{Alan Turing \\ Grace Hopper \\ John Von Neumann \\ Ada Lovelace \\ Charles Babbage}
\usepackage{fancyhdr}\pagestyle{fancy}\usepackage{lastpage}
\lhead{}\chead{}\rhead{}\lfoot{\footnotesize \teamname \ - \semester}\cfoot{}
\rfoot{\footnotesize page \thepage\ of \pageref{LastPage}}
\renewcommand{\headrulewidth}{0.0pt}\renewcommand{\footrulewidth}{0.4pt}
\usepackage[runin]{abstract}\abslabeldelim{\quad}\setlength{\abstitleskip}{-10pt}
\renewcommand{\abstractname}{}\renewcommand{\abstracttextfont}{\color{accentcolor} \small \slshape}

\begin{document}
{\centering \huge \color{accentcolor} \sc \textbf{\department \\ \university} \par}
\vspace{1 in}
{\centering \huge \color{accentcolor} \sc \textbf{\docname \\ \coursename \\ \semester} \par}
\vspace{0.5 in}
\begin{figure}[h!]\centering\includegraphics[width=0.60\textwidth]{images/test_image}\end{figure}
\vspace{0.5 in}
{\centering \huge \color{accentcolor} \sc \textbf{\teamname \\ \productname} \par}
\vspace{0.5 in}{\centering \large \sc \textbf{\authors} \par}
\newpage
\begin{versionhistory}
  \vhEntry{0.1}{09.02.2026}{GH}{document creation}
  \vhEntry{1.0}{09.15.2026}{AT|GH|CB}{official release}
\end{versionhistory}
\newpage\tableofcontents\newpage

\section{Problem Statement}
A county records office receives thousands of paper forms each month containing
handwritten numeric fields (meter readings, account numbers, dated amounts). Staff
key these into a database by hand --- slow, costly, and error-prone, with mistakes
surfacing only later as billing disputes. The office needs a faster, more consistent
way to turn handwritten numeric fields on scanned forms into verified digital data.

\section{Methodology}
Team Perceptron will build \productname, a service that ingests scanned form images,
locates the numeric fields, recognizes the handwritten digits with a machine-learning
model, and returns structured, confidence-scored data. Low-confidence fields are
routed to a human-in-the-loop review screen rather than accepted blindly. In short:
we are building software that reads handwritten numbers off scanned forms and hands
back clean data plus a short review queue.

\section{Value Proposition}
For the sponsor, DigitDesk cuts manual keying effort and transcription errors while
keeping a human check where the model is unsure, lowering downstream dispute costs.
For the university, the project exercises a complete, responsible ML pipeline ---
data handling, training, evaluation, calibration, and human oversight --- with a
public benchmark (MNIST) available for early development.

\section{Development Milestones}
\begin{itemize}
  \item Project Charter first draft - September 2026
  \item System Requirements Specification - October 2026
  \item Architectural Design Specification - November 2026
  \item Demonstration of digit classifier on MNIST - December 2026
  \item Detailed Design Specification - February 2027
  \item Demonstration of field segmentation on real forms - March 2027
  \item Demonstration of review workflow \& confidence routing - April 2027
  \item CoE Innovation Day poster presentation - April 2027
  \item Final Project Demonstration (end-to-end digitization) - May 2027
\end{itemize}
\newpage

\section{Background}
The office currently double-keys critical fields to catch errors, doubling labor.
They approached the CSE program for an ML-assisted alternative that keeps a human in
the loop for auditability. No prior relationship; requirements gathered via two
interviews and a sample of redacted forms.

\section{Related Work}
Handwritten-digit recognition is a classic ML benchmark (MNIST) \cite{Rubin2012},
and commercial OCR/IDP products exist but are priced per-page and often cloud-only,
which conflicts with the office's data-handling rules. DigitDesk targets an
on-premises, auditable, confidence-routed pipeline. (Add $\geq$5 IEEE-format
references in the final version.)

\section{System Overview}
Scanned images enter a preprocessing stage (deskew, denoise, field segmentation);
segmented digit crops go to the recognition model, which emits digits with
confidence; results above a threshold are stored, and the rest enter a review queue
served through a small web UI. See the ADS for the layered design.

\section{Project Management Approach}
Iterative Scrum \cite{Rubin2012, SchwaberSutherland2020} per ISO/IEC/IEEE 16326:2019
\cite{ISO16326} within ISO/IEC/IEEE 12207:2017 \cite{ISO12207}. Two-week sprints;
estimation in task-hours; progress by burndown and by a tracked model-quality metric
(field-level accuracy). Data and model versions are tracked alongside code in Git.

\section{Roles \& Responsibilities}
Sponsor contact: records-office supervisor. Product Owner: Grace Hopper. Scrum
Master: Alan Turing. Data/preprocessing: John Von Neumann; Model/eval: Ada Lovelace;
API/storage: Charles Babbage; Review UI: Grace Hopper; Integration/test: Alan Turing.

\section{Cost Proposal}
Primarily cloud/GPU compute for training and an on-prem inference server; no custom
hardware. Funded by the default CSE department stipend. Largest line item is GPU
training hours.

\section{Facilities \& Equipment}
Senior Design lab workstations, department GPU allocation for training, and an
on-prem test server matching the sponsor environment. A sample of redacted real
forms is required by the 3rd sprint.

\section{Assumptions}
\begin{itemize}
  \item Redacted real-form samples will be available by the 3rd sprint.
  \item Scans are provided at $\geq$200\,dpi grayscale.
  \item Field locations are consistent across a given form template.
  \item Public MNIST data may be used for early model development.
  \item The sponsor can label a validation set of real fields.
\end{itemize}

\section{Constraints}
\begin{itemize}
  \item Final demonstration must be completed by May 1st, 2027.
  \item Total development cost must not exceed \$800.
  \item All form data must remain on sponsor-approved (on-prem) infrastructure.
  \item Personally identifiable fields must be redacted before any team access.
  \item Any accepted-without-review field must meet the agreed accuracy target.
\end{itemize}

\section{Risks}
\begin{table}[h]\resizebox{\textwidth}{!}{
\begin{tabular}{|l|l|l|l|}
\hline
\textbf{Risk description} & \textbf{Probability} & \textbf{Loss (days)} & \textbf{Exposure (days)} \\ \hline
Real handwriting harder than MNIST (domain gap) & 0.55 & 15 & 8.3 \\ \hline
Field-segmentation fails on skewed/noisy scans & 0.40 & 12 & 4.8 \\ \hline
Insufficient labeled real data & 0.45 & 10 & 4.5 \\ \hline
On-prem GPU/server access delayed & 0.25 & 12 & 3.0 \\ \hline
Confidence calibration poor (bad routing) & 0.30 & 8 & 2.4 \\ \hline
\end{tabular}}
\caption{Highest-exposure project risks}
\end{table}

\section{Documentation \& Reporting}
Documentation deliverables tracked in Git and released at their planning sprints;
recurring sprint artifacts and engineering notebooks maintained each sprint.
Closeout: the deployed service, model card, poster, web page, demo video, documented
source (MIT), deployment scripts, and an operator/reviewer manual.

\bibliographystyle{reference/IEEEtran_custom}
\bibliography{reference/refs}{}
\end{document}
